\section{Related Work}\label{sec:related}

\textbf{Mathematical tools.} The one-step problem of Section~\ref{sec:setup} is an I-projection of a distribution onto a linear family, whose solution is an exponential tilt of the reference \cite{csiszar1975idivergence}; its Lagrange dual is a convex program \cite{boyd2004convex}. The chain rule for relative entropy used in Section~\ref{sec:theory} is standard \cite{cover2006elements}. The optimal event-level policy is a Doob $h$-transform of the reference process \cite{doob1957conditional}, and the sequential form of KL-regularized control, in which the optimal controlled transition reweights the uncontrolled one by an exponentiated value function, is the theory of linearly solvable control \cite{kappen2005linear,todorov2007linearly}. The change in the probability of an event under a changed policy as an expected sum of per-step advantages is the performance-difference identity of reinforcement learning \cite{kakade2002approximately}. Equalizing the marginal cost of distortion across components is the equal-slope condition of constrained allocation \cite{cover2006elements,boyd2004convex}; in constrained Markov decision processes a cumulative constraint is handled by a single Lagrange multiplier, whereas a constraint imposed at every step is a different and more restrictive class \cite{altman1999constrained,achiam2017constrained}. We do not claim any of these identities as new.

\textbf{Local token control and global distortion.} That constraining or reweighting a language model token by token induces a sequence distribution different from a sequence-level target is established outside the repetition-loop setting. Controlled decoding \cite{mudgal2024controlled} learns a value function to steer per-token choices toward a KL-regularized reward optimum; twisted Sequential Monte Carlo \cite{zhao2024twisted} and its extension to syntactic and semantic constraints \cite{loula2025syntactic} correct a token-level proposal by particle reweighting because the proposal alone does not sample from the target; adaptive weighted rejection sampling \cite{lipkin2025awrs} makes constrained token-level sampling exact for a broad class of constraints. We do not claim that a gap between local and global control is new. The contribution is narrower: for repetition loops, the gap is a statement about holding a budget fixed rather than a price, it has a closed-form size to second order, and it can be computed exactly in a small system and checked on a language model.

\textbf{Decoding strategies for repetition.} Temperature sampling rescales logits by $1/T$. Nucleus sampling \cite{holtzman2020curious} truncates to the smallest set of tokens with cumulative probability above $p$. Repetition penalty \cite{keskar2019ctrl} divides the logit of every previously generated token by a constant. Typical sampling \cite{meister2023typical} and Mirostat \cite{basu2021mirostat} control the surprisal of the sampled token. No-repeat $n$-gram blocking \cite{paulus2018deep} zeroes the probability of any token that would complete an already-seen $n$-gram. History-aware baselines used here are a longest-suffix-match heuristic (SuffixMatch, our own and not attributed to any publication), an FSD-style decoder adapted from Frustratingly Simple Decoding \cite{yang2023frustratingly}, and a decoder that applies the penalty of the LZ Penalty \cite{ginart2025lzpenalty} with adapted history handling; their relation to the published methods is stated in Appendix~\ref{app:baselines}. Loop onset is treated as a survival time, estimated with Kaplan--Meier \cite{kaplan1958nonparametric}; the character-level models are decoder-only Transformers \cite{vaswani2017attention}.
